% !TEX root = ../main.tex
\chapter{What Is Spacetime Engineering?}\label{ch:what-is}

\begin{goals}
\begin{itemize}
\item What it means to engineer spacetime, and how that differs from
  engineering inside a gravitational field that nature has already fixed.
\item The two directions in which Einstein's equation can be read, and why
  a spacetime engineer reads it \enquote{backwards}.
\item The landmark designs of the field (traversable wormholes, warp drives
  and Krasnikov tubes) and what each one is meant to do.
\item Why these designs are hard to build: the unusual matter they need, the
  limits quantum physics places on that matter, and the threat they pose to
  cause and effect.
\item How the rest of the book is organized.
\end{itemize}
\end{goals}

This chapter is a tour. It introduces the ideas of the whole subject with as
little mathematics as possible, so that you know where the later chapters are
heading and why each one matters. Every claim made here in words is made again
later in equations, and most of them you will derive yourself.

\section{You already depend on engineered time}\label{sec:gps}

Every time a phone finds its position by satellite, it trusts clocks that are
orbiting about \SI{20000}{\kilo\metre} overhead. The navigation satellites
broadcast the time on their onboard atomic clocks, and your receiver works out
how far away each satellite is from how long its signal took to arrive. Light
travels about \SI{30}{\centi\metre} in a nanosecond, so the clocks must agree
with one another to within a few tens of nanoseconds for the position to be
good to a few metres.

Here is the remarkable part. Two identical atomic clocks, one on the ground and
one in a navigation satellite, do not tick at the same rate. The satellite
clock gains about 45 microseconds per day because it sits higher in Earth's
gravity, where time runs faster. It loses about 7 microseconds per day because
it moves at almost \SI{4}{\kilo\metre\per\second}, and moving clocks run slow.
The net effect is that the orbiting clock runs fast by roughly 38 microseconds
every day. Left uncorrected, that offset would produce a ranging error of about
\SI{11}{\kilo\metre} after a single day, since light covers that distance in
38 microseconds. The system works because its designers account for both
effects \parencite{Ashby2003}. You will calculate both numbers yourself in
\cref{ch:geometry}.

Navigation engineers are working \emph{within} a geometry that nature has
fixed. The rates at which clocks tick at different altitudes are set by the
mass of the Earth, and the engineers' job is to measure those rates and correct
for them. Spacetime engineering asks the reverse question. Suppose we could
choose the geometry. What pattern of clock rates, distances and light paths
would we want? What would such a spacetime do to travellers, to light, and to
the ordinary matter around it? And what would it take to produce it?

\section{Spacetime has a geometry}\label{sec:spacetime-geometry}

Special relativity teaches that space and time are not separate. Observers in
relative motion disagree about lengths and durations. They agree, however, on a
combination of the two called the \emph{interval} between events, and on the
speed of light. Together, space and time form a four-dimensional
\emph{spacetime}, and the interval plays the role that distance plays in
ordinary geometry.

General relativity takes one further step. It says that spacetime is not a
fixed, flat stage. Its geometry can be curved, and gravity \emph{is} that
curvature. An apple falls because it follows the straightest possible path
through a curved spacetime, and clocks at different heights tick at different
rates because the geometry of time itself varies from place to place. The
physicist John Wheeler summarized the theory in one sentence: spacetime tells
matter how to move, and matter tells spacetime how to curve.

The rule for the second half of that sentence is \emph{Einstein's equation}
\parencite{Einstein1916}. In words it says
\begin{equation}
  \bigl(\text{curvature of spacetime}\bigr)
  = \frac{8\pi G}{c^{4}}\,
  \bigl(\text{energy, momentum, pressure and stress of matter}\bigr).
  \label{eq:einstein-words}
\end{equation}
The left side is built from the geometry. The right side describes everything
that has energy: matter, light, fields, and even the pressure and tension
inside a material. \Cref{ch:matter-geometry} turns both sides into precise
mathematical objects, and by the end of that chapter you will be able to use
the equation yourself.

The constant in front deserves attention right away, because it sets the scale
of the entire subject. In SI units,
\begin{equation}
  \frac{c^{4}}{G} \approx \SI{1.2e44}{\newton}.
  \label{eq:stiffness}
\end{equation}
This is a force, and an enormous one. Roughly speaking, bending spacetime so
that it curves appreciably over a distance $L$ requires stresses of order
$c^{4}/(G L^{2})$. For $L = \SI{1}{\metre}$ that is about
$\SI{e44}{\pascal}$, more than thirty orders of magnitude beyond the strength of
the strongest materials we know. Spacetime is extraordinarily stiff. This single
number explains why the geometry around you is so nearly flat, and it is the
first reason spacetime engineering is hard.

\begin{checkpoint}
The Earth's gravity bends spacetime enough to make clocks on the ground and in
orbit disagree measurably. How can that be, if spacetime is so stiff? What
property of the Earth makes up for the huge constant in
\cref{eq:stiffness}?
\end{checkpoint}

\section{Reading Einstein's equation in both directions}\label{sec:both-directions}

Physicists usually read Einstein's equation from right to left. They describe
the matter first (a star, a galaxy, the contents of the universe) and then
solve the equation for the geometry that matter produces. This is how the
geometry around the Sun was found, and the geometry of black holes, and the
expanding universe of cosmology. It is also hard: the equation is a set of
coupled, nonlinear partial differential equations, and only highly symmetric
problems can be solved by hand.

A spacetime engineer can read the equation from left to right instead. Start by
deciding what you want the spacetime to do: carry a traveller quickly between
two points, keep the traveller's clock in step with clocks at home, protect the
traveller from tidal forces. Write down a geometry that does those things. Then
compute its curvature, and read off from \cref{eq:einstein-words} the matter
that must be present to produce it. We will call that matter the
\emph{demanded} matter of the geometry. In this direction the mathematics is
easy: computing curvature from a given geometry takes only differentiation.

Michael Morris and Kip Thorne introduced this way of working in 1988, in a paper
written for students \parencite{Morris1988a}. The literature now calls it
\emph{metric engineering}, because the geometry is specified by its metric,
the object you will meet in \cref{ch:geometry}
\parencite{McMonigal2012,Schuster2023}. Morris and Thorne designed spacetimes
containing traversable wormholes by first listing the properties a traveller
would need, then choosing a geometry with those properties, and only then
asking what matter the geometry would require.

The catch comes at the last step. Einstein's equation always returns an answer:
\emph{every} geometry has some demanded matter. Nothing guarantees that such
matter exists. The demand might call for energy densities that are negative, or
tensions larger than any material can bear, or both. The real work of spacetime
engineering begins there, and it proceeds in a loop (\cref{fig:design-loop}).

\begin{figure}[htbp]
\centering
\begin{tikzpicture}[
  node distance=0.55cm and 0.9cm,
  stage/.style={draw=goalblue, thick, rounded corners=3pt, fill=goalblue!6,
    align=center, text width=5.6cm, minimum height=1.05cm, font=\small},
  arrow/.style={-{Stealth[length=2.5mm]}, thick, goalblue},
]
\node[stage] (want) {What should the spacetime do?\\\footnotesize carry a traveller, shorten a trip, keep clocks in step};
\node[stage, below=of want] (geom) {Choose a geometry that does it};
\node[stage, below=of geom] (demand) {Compute the matter it demands};
\node[stage, below=of demand] (test) {Test the demand against physics\\\footnotesize energy conditions, quantum limits, causality, stability, effects on travellers};
\node[stage, below=of test] (source) {Find matter and fields that could supply it};
\draw[arrow] (want) -- (geom);
\draw[arrow] (geom) -- (demand);
\draw[arrow] (demand) -- (test);
\draw[arrow] (test) -- (source);
\draw[arrow] (test.east) -- ++(1.1,0) |- (geom.east);
\draw[arrow] (source.east) -- ++(1.5,0) |- node[pos=0.25, right=3pt, font=\footnotesize, align=left, text=goalblue] {revise the\\geometry} (geom.east);
\end{tikzpicture}
\caption{The design loop of spacetime engineering. A geometry is chosen for what
it does, and its demanded matter is tested against physics. When no known
matter can supply the demand, the geometry is revised. Most of this book
teaches the steps inside this loop.}
\label{fig:design-loop}
\end{figure}

The loop can also run the other way. Some recent work starts from a matter
distribution with sensible properties, such as a shell of ordinary matter,
solves for the geometry it produces, and then asks what that geometry can do.
Both directions appear in this book. The geometry-first direction is the more
common one, and it is the one that most clearly exposes what a design costs.

\begin{checkpoint}
Why is it always possible to compute the demanded matter of a geometry, while
solving for the geometry produced by given matter is often impossible to do by
hand?
\end{checkpoint}

\section{A gallery of engineered spacetimes}\label{sec:gallery}

Three designs define the field. Each was invented to answer the question
\enquote{could we travel between the stars much faster than light takes?}, and
each answers it in a different way.

\subsection{Traversable wormholes}

A wormhole is a tunnel through space: two openings, perhaps light-years apart,
joined by a short passage. \Cref{fig:wormhole} shows the usual way of picturing
one. Take a two-dimensional slice through space around the wormhole, and draw
it as a curved surface. Far from the wormhole the surface is flat. Near it, the
surface curves down into a funnel, narrows to a smallest circle called the
\emph{throat}, and opens out again into a second region of space.

\begin{figure}[htbp]
\centering
\begin{tikzpicture}[scale=0.95]
  % profile curves r(z) = cosh(z)
  \draw[thick, goalblue, domain=-1.9:1.9, samples=60, smooth, variable=\z]
    plot ({cosh(\z)}, {\z});
  \draw[thick, goalblue, domain=-1.9:1.9, samples=60, smooth, variable=\z]
    plot ({-cosh(\z)}, {\z});
  % rims and throat
  \draw[thick, goalblue] (0,1.9) ellipse ({cosh(1.9)} and 0.42);
  \draw[goalblue, dashed] (0,-1.9) ellipse ({cosh(1.9)} and 0.42);
  \draw[goalblue, thick, domain=180:360, samples=40, variable=\t]
    plot ({cosh(-1.9)*cos(\t)}, {-1.9+0.42*sin(\t)});
  \draw[keygold, very thick] (0,0) ellipse (1 and 0.22);
  \draw[goalblue!60] (0,1.0) ellipse ({cosh(1.0)} and 0.3);
  \draw[goalblue!60] (0,-1.0) ellipse ({cosh(1.0)} and 0.3);
  % labels
  \node[keygold!80!black, font=\small, anchor=west] at (1.25,0) {throat};
  \node[font=\small, align=center] at (0,2.75) {opening into region A};
  \node[font=\small, align=center] at (0,-2.75) {opening into region B};
  \draw[-{Stealth}, examplegreen, thick] (-0.25,1.7) .. controls (-0.15,0.5) and (0.15,-0.5) .. (0.25,-1.7);
  \draw[examplegreen!70, thin] (-0.2,1.15) -- (-2.25,1.5);
  \node[examplegreen, font=\footnotesize, anchor=east] at (-2.3,1.5) {traveller};
\end{tikzpicture}
\caption{An embedding diagram of a wormhole. A two-dimensional slice of space
is drawn as a surface. Circles around the wormhole shrink toward the throat and
grow again on the far side: the throat \emph{flares out} in both directions.}
\label{fig:wormhole}
\end{figure}

Wormholes appear as mathematical solutions in general relativity going back to
the 1930s, but the early examples could not be crossed. Morris and Thorne asked
what a wormhole would need in order to be \emph{traversable}: no horizon at the
throat, tidal forces a person could survive, and a crossing that takes a
reasonable time. They found geometries that meet all these requirements, and
then computed what those geometries demand.

\begin{history}{a novel and a wormhole}
Morris and Thorne were prompted to study traversable wormholes in the summer of
1985, when the astronomer Carl Sagan sent Thorne a draft of his novel
\emph{Contact} and asked for help making its gravitational physics as accurate
as possible \parencite{Morris1988a}. Sagan revised the novel's travel method
in light of their preliminary results. The resulting paper, subtitled \enquote{a
tool for teaching general relativity}, launched the modern study of engineered
spacetimes.
\end{history}

The difficulty lies at the throat. A bundle of light rays heading into a
wormhole converges as it approaches the throat. To come out the other side, the
bundle must spread apart again. Ordinary matter, however, always makes light
rays converge: gravity attracts, and it bends light the way a converging lens
does. A throat has to act as a \emph{diverging} lens. Morris and Thorne showed
that this requires the material at the throat to carry a tension greater than
its energy density. Observers passing through the throat at speeds close to
that of light would then measure a \emph{negative} energy density. Matter with
this property is called \emph{exotic}. You will derive the flaring-out
condition and its consequence for yourself in Part~II.

\subsection{Warp drives}

In 1994 Miguel Alcubierre took a different approach \parencite{Alcubierre1994}.
Instead of connecting two distant places by a tunnel, his design moves a region
of space itself. Space contracts in front of a bubble and expands behind it,
and a region of perfectly flat spacetime inside the bubble is carried along
(\cref{fig:warp}). A ship inside that region does not move through its
surroundings at all; it floats in free fall while the bubble carries it.

\begin{figure}[htbp]
\centering
\begin{tikzpicture}[scale=0.9]
  \draw[-{Stealth}, thick, gray] (-5.6,0) -- (5.8,0) node[right, font=\small, text=black] {direction of travel};
  \draw[thick, deeperviolet, dashed] (0,0) circle (1.55);
  \fill[goalblue!15] (0,0) circle (1.05);
  \node[font=\footnotesize, align=center] at (0,0.55) {flat region};
  \draw[fill=goalblue!60, draw=goalblue] (-0.45,-0.14) rectangle (0.45,0.14);
  \node[font=\footnotesize] at (0,-0.5) {ship};
  \foreach \y in {-1.1,-0.45,0.45,1.1} {
    \draw[-{Stealth}, thick, keygold!80!black] (3.9,\y) -- (2.2,\y);
    \draw[-{Stealth}, thick, examplegreen] (-2.2,\y) -- (-3.9,\y);
  }
  \node[font=\small, align=center, text=keygold!70!black] at (3.2,1.75) {space contracts\\ahead};
  \node[font=\small, align=center, text=examplegreen] at (-3.2,1.75) {space expands\\behind};
  \node[font=\footnotesize, text=deeperviolet] at (0,-1.95) {bubble wall};
\end{tikzpicture}
\caption{A schematic warp drive. Space contracts in front of the bubble and
expands behind it, and the flat region inside is carried along. The ship inside
is in free fall.}
\label{fig:warp}
\end{figure}

Alcubierre's bubble has several striking properties. Nothing inside it ever
moves faster than light relative to its own surroundings, yet the bubble as a
whole can move at any speed relative to distant observers. The ship feels no
acceleration. In Alcubierre's design its clock ticks at the same rate as clocks
far away, so the travellers age no more than the people they left behind
during the trip. José Natário later showed that the expansion and contraction
are not essential: a warp drive can carry its passengers with no change in
volume at all \parencite{Natario2002}. What \emph{is} essential is the demand.
The bubble's wall requires negative energy density, just as the wormhole's
throat does.

\subsection{Krasnikov tubes}

Sergei Krasnikov pointed out a problem with the warp drive as a means of
travel: the people inside a bubble moving faster than light cannot send signals
to its front edge, so they cannot create the bubble or steer it
\parencite{Krasnikov1998}. His alternative has the traveller build the
structure on the way out. A ship travels to a distant star at less than light
speed, and as it goes it modifies the spacetime along its path, leaving behind
a tube. The modified geometry inside the tube lets the return trip arrive home
shortly after departure. Allen Everett and Thomas Roman analysed these tubes
further and showed that two of them, placed suitably, could form a time
machine \parencite{Everett1997}.

\subsection{Recent work}

In 2021 several authors proposed warp drives that appeared to require only
positive energy. Examining them closely, Jessica Santiago, Sebastian Schuster
and Matt Visser showed that the positive energy was measured by one special
family of observers only. Other observers see negative energy, and every
physically reasonable warp drive violates the conditions that ordinary matter
satisfies \parencite{Santiago2022}. That episode taught the field a lesson you
will meet again and again in this book: the demand of a geometry must be
checked from the point of view of \emph{every} observer. Current research
includes shells of ordinary matter that carry a flat interior at speeds below
that of light, computer simulations of warp spacetimes, and the search for
fields that could supply negative energy.

\section{The central obstacle}\label{sec:obstacle}

All the designs in the gallery share one requirement. Somewhere, they need
matter that bends light the \enquote{wrong} way.

\subsection{Energy conditions}

Physicists summarize what ordinary matter can do with a small set of
inequalities called \emph{energy conditions}. The simplest to state is the
\emph{weak energy condition}: every observer, however they move, measures a
non-negative energy density. The most important for this book is the
\emph{null energy condition}, which is weaker still. It concerns light rays
passing through matter, and it guarantees that matter always focuses a
bundle of light rays rather than spreading it. Everything you have ever
touched satisfies these conditions, and so do light, electric and magnetic
fields, and all ordinary fluids.

The null energy condition is exactly what shortcuts need to violate. A wormhole
throat must defocus light, and a warp bubble's wall must do the same. In
Chapter~5 you will see that this is no accident: under quite general
assumptions, a spacetime that lets travellers outrun light through an otherwise
ordinary region of space must violate the null energy condition somewhere.

\begin{keyidea}
Designing a shortcut through spacetime is easy on paper. Supplying the matter
it demands is the hard part, because shortcuts demand matter that violates the
null energy condition, and no ordinary matter does.
\end{keyidea}

\subsection{Negative energy in the laboratory}

Negative energy density is not pure fiction. In 1948 Hendrik Casimir predicted
that two parallel, uncharged metal plates in vacuum attract each other
\parencite{Casimir1948}. The plates restrict which electromagnetic vacuum
fluctuations fit between them, and the energy density between the plates is
lower than that of empty space: in this sense it is negative. Steve Lamoreaux
measured the Casimir force in 1997 \parencite{Lamoreaux1997}, and it has since
been measured many times. Quantum physics therefore allows negative energy.

It allows it only in limited amounts. Lawrence Ford and Thomas Roman showed that
quantum fields obey \emph{quantum inequalities}: the more negative the energy
density in a region, the shorter the time it can last and the smaller the region
it can occupy \parencite{Ford1995,Ford1997}. Applied to a warp bubble, these
inequalities force the bubble's wall to be extraordinarily thin. Michael
Pfenning and Ford estimated that a bubble with a radius of \SI{100}{\metre} would
then need about $\SI{6e62}{\kilogram}$ of negative energy for every unit of its
speed in units of $c$ \parencite{Pfenning1997}. That is vastly more than the
total mass of the observable universe. Clever reshaping of the bubble can cut
the requirement dramatically: Chris Van Den Broeck brought it down to a few
solar masses \parencite{VanDenBroeck1999}. That is still an astronomical
amount, and his design bends space on scales of about ten Planck lengths,
far below anything that has ever been probed.

\begin{checkpoint}
The Casimir effect shows that negative energy density exists. Why does that not
settle the question of whether a warp drive can be built?
\end{checkpoint}

\subsection{Causality and control}

Two further obstacles have nothing to do with the amount of energy. The first
is causality. Relativity allows observers in relative motion to disagree about
which of two distant events happens first. If a signal or a traveller could
outrun light, some observers would see it arrive before it left. With two
faster-than-light structures, a traveller could return before departing and
change the past. Stephen Hawking conjectured that physics always prevents
this, a proposal he called the \emph{chronology protection conjecture}
\parencite{Hawking1992}. It remains a conjecture, and any practical design
must show that it creates no such loops.

The second is control. As Krasnikov noticed, the occupants of a warp bubble
moving faster than light are cut off from the bubble's front edge. They can
neither build the bubble nor steer it from inside. Designs that carry
travellers faster than light must therefore be laid out in advance, or driven
by something outside the travellers' reach.

\section{What spacetime engineering is today}\label{sec:today}

With all these obstacles, what is the subject for? Nobody knows how to build
any of the spacetimes in this book, and some of them may be forbidden outright.
Spacetime engineering today is a theoretical discipline, and it earns its place
in three ways.

It teaches general relativity. Designing a spacetime forces you to understand
how every piece of its geometry affects clocks, rulers, light and matter. That
was Morris and Thorne's reason for writing their paper, and it is one reason for
this book.

It maps the boundary of the possible. Some of the obstacles above are theorems
with precise assumptions. Some are estimates that depend on how a design is
drawn. Some are conjectures. Knowing which is which tells you what would have to
change for a design to become possible, and a large part of this book is about
learning to tell them apart.

It develops methods that carry over to other problems: how to design under
severe physical constraints, how to check a long calculation, and how to judge
how far a result can be trusted. Those methods are useful long before anyone
builds a warp drive.

\section{How this book is organized}\label{sec:organization}

The rest of Part~I gives you the tools. \Cref{ch:geometry} develops the geometry
of spacetime: intervals, proper time, the metric, curvature and free fall.
\Cref{ch:matter-geometry} describes matter by its energy, momentum and stress,
states Einstein's equation, and shows how to compute the matter a geometry
demands. The chapters that follow split spacetime into space and time the way
an engineer describes a design, state the energy conditions precisely, explain
what quantum physics allows, and treat causality and the shape of space as a
whole.

Part~II studies the classic designs in detail. Part~III is about design: how
the rate of clocks, the flow of space and the shape of space each affect what a
spacetime demands and what it does to the people inside it. Part~IV asks what
kinds of matter and fields could supply a demand. Part~V takes up dynamics,
the checking of calculations, and the open problems of the field.

In this chapter we have used SI units. From \cref{ch:geometry} onward we
measure time and distance in the same units, which sets $c = 1$, and from
\cref{ch:matter-geometry} onward we also set $G = 1$. Each chapter shows how to
convert back to everyday units whenever a number matters.

\begin{chaptersummary}
\item Spacetime has a geometry that matter bends; gravity is that bending.
  Einstein's equation links the curvature of spacetime to the energy,
  momentum, pressure and stress of matter.
\item The constant $c^{4}/G \approx \SI{1.2e44}{\newton}$ makes spacetime
  extremely stiff. Large changes in geometry require enormous amounts of
  energy and stress.
\item Spacetime engineering reads Einstein's equation backwards: choose a
  geometry for what it does, then compute the matter it demands. Every
  geometry has a demanded matter content; whether that matter exists is the
  real question.
\item The classic designs are traversable wormholes (tunnels through space),
  warp drives (bubbles of flat space carried along by a flow of space) and
  Krasnikov tubes (structures laid down on an outbound trip).
\item All of them demand matter that violates the null energy condition,
  which every ordinary material satisfies. Quantum physics allows negative
  energy density, as the Casimir effect shows, but quantum inequalities limit
  how much and for how long.
\item Faster-than-light designs also face causality (the risk of time
  machines) and control (occupants cannot steer a faster-than-light bubble
  from inside).
\end{chaptersummary}

\begin{checkanswers}
\item[1.1] The stiffness constant multiplies the \emph{density} of energy.
  The Earth's mass, about $\SI{6e24}{\kilogram}$, is concentrated enough that
  its gravitational potential at the surface is about $7\times10^{-10}$ in
  units of $c^{2}$. That tiny fraction is exactly the size of the clock-rate
  differences GPS must correct. Large masses produce small curvature, and
  precise clocks can measure small curvature.
\item[1.2] Computing the demanded matter only requires differentiating the
  geometry, which can always be done. Solving for the geometry means solving
  nonlinear partial differential equations for unknown functions, which is
  possible by hand only in highly symmetric cases.
\item[1.3] The quantum inequalities limit how large and long-lived a region of
  negative energy can be. A warp drive needs large amounts of negative energy
  held in place for the duration of the trip, far beyond what the
  inequalities allow for known quantum fields.
\end{checkanswers}

\begin{problems}
\item \diff{1} \textbf{Two clocks.} A clock sits at the top of a tower of height
  $h = \SI{100}{\metre}$. Near the Earth's surface, a clock higher by $h$ runs
  faster by the fraction $gh/c^{2}$, where $g = \SI{9.8}{\metre\per\second\squared}$.
  How many nanoseconds does the top clock gain over a clock at the base in one
  year?
\item \diff{1} \textbf{The stiffness of spacetime.} Compute $c^{4}/G$ in newtons
  from $c = \SI{2.998e8}{\metre\per\second}$ and
  $G = \SI{6.674e-11}{\newton\metre\squared\per\kilogram\squared}$. Then
  estimate the stress, in pascals, needed to curve spacetime appreciably over a
  distance of \SI{1}{\kilo\metre}, and compare it with the tensile strength of
  steel, about $\SI{e9}{\pascal}$.
\item \diff{2} \textbf{A day of GPS error.} Using the numbers in
  \cref{sec:gps}, find how long it takes for an uncorrected satellite clock to
  drift by the \SI{30}{\nano\second} that corresponds to about \SI{10}{\metre}
  of position error.
\item \diff{2} \textbf{Converging and diverging lenses.} Explain in your own
  words why a wormhole throat must defocus light, and why ordinary matter
  cannot do that. (Hint: think about how a bundle of light rays behaves as it
  enters the throat, passes through, and leaves on the other side.)
\item \diff{2} \textbf{How much is too much?} Express the Pfenning–Ford estimate
  of $\SI{6e62}{\kilogram}$ in solar masses, taking
  $\Msun = \SI{2.0e30}{\kilogram}$. Van Den Broeck's design needs a few solar
  masses. By what factor did reshaping the bubble reduce the requirement?
\item \diff{3} \textbf{Which assumptions?} The quantum inequalities rest on
  assumptions: the kind of quantum field, the flatness of spacetime over the
  sampling time, and the way the energy is averaged. For each, describe how a
  designer might try to take advantage of it, and what you would check to see
  whether the attempt succeeds. Return to your answer after
  Chapter~6.
\end{problems}

\begin{furtherreading}
\item[\textcite{Morris1988a}] The paper that started the field, written for
  students and still one of the best introductions to engineered spacetimes.
\item[\textcite{Alcubierre1994}] A five-page paper introducing the warp drive,
  readable after \cref{ch:matter-geometry}.
\item[\textcite{Everett2011}] A book-length, non-technical account of warp
  drives, wormholes and time travel by two researchers who helped shape the
  field.
\item[\textcite{Visser1995b}] The standard monograph on wormholes, for readers
  who want the full mathematical treatment.
\item[\textcite{Kontou2020}] A modern review of energy conditions in classical
  and quantum physics, useful as a reference throughout this book.
\item[\textcite{Ashby2003}] The complete relativistic treatment of the Global
  Positioning System.
\end{furtherreading}
